\scanpage{111}
% Archive foliation: 112, inverted at upper right. Reused unrelated typescript; printed page 15, n°382.
\begin{center}-- 15 --\hfill n\textsuperscript{o} 382\end{center}
qu'il existe un nombre réel $M$ tel que
\[
\lVert b_s\rVert\leq M^{|s|}\qquad\text{pour tout }s\ne0.
\]
Les séries absolument convergentes forment un sous-espace vectoriel $E\{\{X\}\}$ de $E[[X]]$ ; c'est un module sur l'algèbre $K\{\{X\}\}$ des séries convergentes scalaires. Lorsque $E$ est un espace vectoriel de dimension finie, l'application canonique de $E\otimes_K K\{\{X\}\}$ dans $E\{\{X\}\}$ est un isomorphisme.

3.3.3. Soit $f=\sum b_sX^s$ une série convergente à coefficients dans $E$. Soit $J(f)$ le sous-ensemble de $\underline{\underline{R}}_+^n$ formé des points $r=(r_1,\ldots,r_n)$, $r_i>0$, tels que la série à termes positifs $\sum\lVert b_s\rVert r^s$ soit convergente. Si $r\in J(f)$, la série $f$ est absolument convergente en tout point du polydisque $B(r)$. L'intérieur $I(f)$ de $J(f)$ s'appelle \uline{l'indicatrice de convergence} de $f$. Pour qu'un point $r=(r_1,\ldots,r_n)$, $r_i>0$, appartienne à $I(f)$, il faut et il suffit qu'il existe un point $r'=(r'_1,\ldots,r'_n)$, avec $r'_i>r_i$, tel que
\[
\operatorname{Sup}\lVert b_s\rVert r'^{s}<+\infty.
\]
Soit $\Omega(f)$ le sous-ensemble de $\underline{\underline{R}}^n$ formé des points
\[
\log r=(\log r_1,\ldots,\log r_n),
\]
où $r$ parcourt $I(f)$. C'est un sous-ensemble ouvert \uline{convexe} de $\underline{\underline{R}}^n$ ; de plus, il existe un nombre réel $k$ (dépendant de $f$) tel que $\Omega(f)$ contienne tous les points $(a_1,\ldots,a_n)$ tels que $a_i\leq k$.

3.3.4. L'ensemble $D(f)=\bigcup_{r\in I(f)}B(r)=\bigcup_{r\in J(f)}B(r)$ est un ouvert non vide de $K^n$ ; il est contenu dans l'intérieur $\overset{\circ}{C}(f)$ de l'ensemble $C(f)$ des points où $f$ converge absolument ; lorsque $K$ est égal à $\underline{\underline{R}}$ ou à $\underline{\underline{C}}$, on a $D(f)=\overset{\circ}{C}(f)$. L'ouvert $D(f)$ est appelé le \uline{domaine de convergence stricte} de $f$.

3.4 -- \uline{Opérations sur les séries convergentes}.

3.4.1. Soit $E'$ un espace vectoriel normé complet et soit $u:E\longrightarrow E'$ une application linéaire continue. Soit $f=\sum b_sX^s$ une série convergente à coefficients dans $E$. Alors la série
